\pretitle{Agradecimentos}

% =========================================================
% AGRADECIMENTOS
% FINALIDADE: reconhecer pessoas e instituições que contribuíram de
%   maneira relevante para o trabalho.
% ESCREVA: um ou mais parágrafos, no lugar dos parágrafos de orientação
%   abaixo. O título "Agradecimentos" é automático.
% OPCIONAL: para retirar, use \IncluirAgradecimentos{nao} em
%   1-dados/dados.tex.
% CUIDADO: com bolsa ou financiamento, use a redação exigida pela agência.
%   Bolsistas CAPES (Portaria CAPES nº 206/2018) devem incluir:
%   O presente trabalho foi realizado com apoio da Coordenação de
%   Aperfeiçoamento de Pessoal de Nível Superior -- Brasil (CAPES) --
%   Código de Financiamento 001.
%   Outras agências de fomento possuem exigências próprias.
% =========================================================

Elemento opcional. Texto em que o(a) autor(a) registra agradecimentos às pessoas e às instituições que contribuíram de maneira relevante para a elaboração do trabalho.

Quando a pesquisa tiver contado com bolsa ou financiamento, o agradecimento à agência de fomento deve seguir a redação exigida pela instituição financiadora.
